\begin{proof}[Proof of \cref{obj:lem:released-arm-cell-masses}]
The law \(P\) is a probability law by \cref{obj:def:causal-law-class}. Its score marginal is \(H\) by \cref{obj:ass:score-marginal}, and its released-record marginal is \(P_{\mathrm{rel}}\). Both marginals are therefore probability laws. % lean: CausalSmith.PartialID.UnlinkedPropensityAte.compatibleCellMasses_of_randomizedCausalLaw

Fix \(r\in\mathcal R_J\). By \cref{obj:ass:deterministic-release}, \(R=g(E)\) almost surely. Conditioning on \((E,Y_0,Y_1)\) and applying \cref{obj:ass:randomized-assignment} gives
\[
\begin{aligned}
P_{\mathrm{rel}}(A=1,R=r)
&=P(A=1,g(E)=r)\\
&=\mathbb E_P\!\left[E\,\mathbf 1_{\{g(E)=r\}}\right]\\
&=\int_{C_r}e\,H(de).
\end{aligned}
\] % lean: CausalSmith.PartialID.UnlinkedPropensityAte.compatibleCellMasses_of_randomizedCausalLaw

For the control arm, partition the event \(E\in C_r\) by treatment assignment. The score marginal condition and the preceding identity yield
\[
\begin{aligned}
P_{\mathrm{rel}}(A=0,R=r)
&=P(E\in C_r)-P(A=1,E\in C_r)\\
&=H(C_r)-\int_{C_r}e\,H(de)\\
&=\int_{C_r}(1-e)\,H(de)
=q_{0r}.
\end{aligned}
\] % lean: CausalSmith.PartialID.UnlinkedPropensityAte.compatibleCellMasses_of_randomizedCausalLaw
The integral subtraction is valid because the score is bounded on \(\mathcal E\). Finally, \(p_1(e)=e\) in \cref{obj:def:arm-cell-primitives}, so the treated-arm identity is \(P_{\mathrm{rel}}(A=1,R=r)=q_{1r}\). These two identities cover every \(a\in\mathcal A\). % lean: CausalSmith.PartialID.UnlinkedPropensityAte.compatibleCellMasses_of_randomizedCausalLaw
\end{proof}
